
\documentclass[11pt]{article}

% --------------------------------------------------
% Packages
% --------------------------------------------------
\usepackage[margin=1in]{geometry}
\usepackage{amsmath,amssymb,amsthm,mathtools}
\usepackage{hyperref}

% --------------------------------------------------
% Formatting
% --------------------------------------------------
\setlength{\parindent}{0pt}
\setlength{\parskip}{0.3cm}

% --------------------------------------------------
% Theorem Environments
% --------------------------------------------------
\newtheorem{theorem}{Theorem}[section]
\newtheorem{lemma}[theorem]{Lemma}
\newtheorem{proposition}[theorem]{Proposition}
\newtheorem{corollary}[theorem]{Corollary}

\theoremstyle{definition}
\newtheorem{definition}[theorem]{Definition}
\newtheorem{example}[theorem]{Example}

\theoremstyle{remark}
\newtheorem{remark}[theorem]{Remark}

% --------------------------------------------------
% Title
% --------------------------------------------------
\title{\textbf{Congruence and Sylvester's Law of Inertia}\\
\large An Elementary Exposition with Applications In Mathematical Statistics}

\author{Biajid Chowdhury
}
\date{\today}

\begin{document}

\maketitle

% --------------------------------------------------
% Abstract
% --------------------------------------------------
\begin{abstract}

In many modern undergraduate linear algebra curricula,
the concept of \emph{inertia} receives relatively little
attention, if it is discussed at all. Consequently,
students may become familiar with the diagonalization
of symmetric matrices and the notion of matrix
similarity without fully appreciating the distinction
between similarity and congruence. This distinction,
although subtle at first glance, is fundamental to
the theory of quadratic forms.

In this note, we aim to clarify these ideas and
develop Sylvester's Law of Inertia from first
principles. Beginning with quadratic forms and
changes of basis, we introduce matrix congruence,
define the inertia of a real symmetric matrix,
and establish the remarkable fact that the numbers
of positive, negative, and zero eigenvalues remain
unchanged under congruence.

Particular attention is devoted to a detailed,
self-contained proof of the theorem, emphasizing
the underlying linear algebra rather than merely
stating the result.

\end{abstract}

\tableofcontents
\newpage

% --------------------------------------------------
% 1. Introduction
% --------------------------------------------------
\section{Introduction}

One of the fundamental results of linear algebra is the
spectral theorem, which guarantees that every real symmetric
matrix can be diagonalized by an orthogonal change of basis.
While diagonalization and matrix similarity are familiar
topics in undergraduate linear algebra, the closely related
notion of \emph{congruence} often receives less attention.

Similarity and congruence arise in different mathematical
settings. The transformation $S^{-1}AS$ describes a change
of basis for a linear operator, whereas $S^TAS$ arises
naturally from a change of coordinates in a quadratic form.
Indeed, if a real symmetric matrix $A$ defines
\[
    q(x)=x^TAx,
\]
then the substitution $x=Sy$, with $S$ invertible, gives
\[
    q(x)=y^T(S^TAS)y.
\]
Thus, $A$ and $S^TAS$ represent the same quadratic form
in different coordinate systems and are said to be
\emph{congruent}.

Unlike similarity, congruence does not generally preserve
individual eigenvalues. Remarkably, however, it preserves
the numbers of positive, negative, and zero eigenvalues.
These three numbers constitute the \emph{inertia} of a
real symmetric matrix, and their invariance under
congruence is known as \emph{Sylvester's Law of Inertia}.

The purpose of this note is to develop this theorem from
first principles. We begin with a brief review of quadratic forms and
changes of basis, distinguish congruence from similarity,
and introduce the notions of inertia, rank, index, and
signature. We then present a detailed proof of Sylvester's
Law of Inertia using elementary tools from linear algebra,
followed by its canonical form and applications in
mathematical statistics and theoretical physics.
We conclude with a brief historical account of James
Joseph Sylvester and the development of the theorem.

\section{Prerequisites}

This note is intended for students and researchers
who have
completed a first course in linear algebra. Readers are
expected to be well versed in eigenvalues and eigenvectors,
eigendecomposition, the spectral theorem for real symmetric
matrices, and the theory of quadratic forms.

Familiarity with changes of basis, orthonormal bases,
matrix rank, null spaces, and the rank-nullity theorem
is also assumed. In particular, the reader should
understand why real symmetric matrices admit orthonormal
eigenbases and how a quadratic form can be represented
by a symmetric matrix.

Although these concepts form the mathematical foundation
of our discussion, the essential ideas will be briefly
revisited whenever necessary. No knowledge of Sylvester's
Law of Inertia or the theory of matrix congruence is
assumed.

\section{Quadratic Forms and Changes of Basis}

Before introducing matrix congruence, we briefly recall
the notion of a quadratic form and examine how its
matrix representation changes under a change of basis.

\subsection{Quadratic Forms}

A quadratic form on $\mathbb{R}^n$ is a scalar-valued
function of the form
\[
    q(x)=x^TAx,
\]
where $A\in\mathbb{R}^{n\times n}$ is a real symmetric
matrix and $x\in\mathbb{R}^n$ is a column vector.

For example, consider
\[
    q(x_1,x_2)=2x_1^2+2x_1x_2+3x_2^2.
\]
This quadratic form can be represented as
\[
    q(x)=
    \begin{pmatrix}x_1&x_2\end{pmatrix}
    \begin{pmatrix}
        2&1\\
        1&3
    \end{pmatrix}
    \begin{pmatrix}x_1\\x_2\end{pmatrix}.
\]

Notice that the coefficient of the mixed term
$x_1x_2$ is distributed equally between the
corresponding off-diagonal entries of the
symmetric matrix.

More generally, every quadratic form on
$\mathbb{R}^n$ admits a unique symmetric
matrix representation relative to a fixed basis.

\subsection{Changing the Basis}

Suppose that the quadratic form
\[
    q(x)=x^TAx
\]
is initially expressed in the standard basis
of $\mathbb{R}^n$.

Let
\[
    \mathcal{B}=\{v_1,v_2,\ldots,v_n\}
\]
be another basis, and define the change-of-basis
matrix
\[
    S=\begin{pmatrix}v_1&v_2&\cdots&v_n\end{pmatrix}.
\]

Since the columns of $S$ form a basis,
$S$ is invertible.

If $y$ denotes the coordinate vector of $x$
relative to the new basis $\mathcal{B}$, then
\[
    x=Sy.
\]

Substituting into the quadratic form gives
\[
\begin{aligned}
    q(x)
    &=x^TAx\\
    &=(Sy)^TA(Sy)\\
    &=y^TS^TASy.
\end{aligned}
\]

Therefore, if we define
\[
    \boxed{B=S^TAS},
\]
and the quadratic form in the new coordinates becomes
\[
    q(Sy)=y^TBy.
\]

Thus, $A$ and $B$ represent the same quadratic
form relative to two different bases.

This motivates the following definition.

\begin{definition}[Matrix Congruence]
Two real symmetric matrices
$A,B\in\mathbb{R}^{n\times n}$
are said to be \emph{congruent} if there exists
an invertible matrix $S$ such that
\[
    \boxed{B=S^TAS}.
\]
\end{definition}

\begin{remark}
A change of basis does not alter the quadratic
form itself; it changes only its matrix
representation. Moreover, the change-of-basis
matrix need not be orthogonal. Any invertible
matrix is permitted.
\end{remark}

\subsection{An Illustrative Example}

Consider the quadratic form
\[
    q(x_1,x_2)=2x_1^2+2x_1x_2+3x_2^2,
\]
whose symmetric matrix representation is
\[
    A=
    \begin{pmatrix}
        2&1\\
        1&3
    \end{pmatrix}.
\]

Suppose we introduce a new basis
\[
    v_1=
    \begin{pmatrix}1\\0\end{pmatrix},
    \qquad
    v_2=
    \begin{pmatrix}1\\1\end{pmatrix}.
\]

The corresponding change-of-basis matrix is
\[
    S=
    \begin{pmatrix}
        1&1\\
        0&1
    \end{pmatrix}.
\]

The matrix of the quadratic form in the new
basis is
\[
\begin{aligned}
    B
    &=S^TAS\\
    &=
    \begin{pmatrix}
        1&0\\
        1&1
    \end{pmatrix}
    \begin{pmatrix}
        2&1\\
        1&3
    \end{pmatrix}
    \begin{pmatrix}
        1&1\\
        0&1
    \end{pmatrix}\\[6pt]
    &=
    \boxed{
    \begin{pmatrix}
        2&3\\
        3&7
    \end{pmatrix}}.
\end{aligned}
\]

To verify this directly, observe that $x=Sy$
implies
\[
    x_1=y_1+y_2,
    \qquad
    x_2=y_2.
\]

Substituting into the original quadratic form,
\[
\begin{aligned}
    q(x)
    &=2(y_1+y_2)^2
      +2(y_1+y_2)y_2+3y_2^2\\
    &=2y_1^2+6y_1y_2+7y_2^2\\
    &=
    \begin{pmatrix}y_1&y_2\end{pmatrix}
     \boxed{
    \begin{pmatrix}
        2&3\\
        3&7
    \end{pmatrix}}.
    \begin{pmatrix}y_1\\y_2\end{pmatrix}.
\end{aligned}
\]

We have therefore obtained two different
symmetric matrices representing exactly the
same quadratic form in different coordinate
systems.

The important question is now the following:
although congruence preserves the quadratic
form, does it also preserve its eigenvalues?

As we shall see in the next section, the
answer is generally no. This distinction
separates congruence from matrix similarity
and ultimately leads to Sylvester's Law
of Inertia.


\section{Similarity Versus Congruence}

Although similarity and congruence are both associated
with changes of basis, they describe fundamentally
different mathematical relationships. It is important
to distinguish them before proceeding to Sylvester's
Law of Inertia.

\subsection{Similarity and Linear Operators}

Two matrices $A,B\in\mathbb{R}^{n\times n}$ are
said to be \emph{similar} if there exists an
invertible matrix $S$ such that
\[
    B=S^{-1}AS.
\]

Similar matrices represent the same linear operator
relative to different bases. More importantly,
they have the same characteristic polynomial:
\[
\begin{aligned}
\det(\lambda I-B)
&=\det(\lambda I-S^{-1}AS)\\
&=\det\bigl(S^{-1}(\lambda I-A)S\bigr)\\
&=\det(\lambda I-A).
\end{aligned}
\]

Consequently, similar matrices possess exactly
the same set of eigenvalues, including their algebraic
multiplicities. In particular, similarity preserves
the trace, determinant, and rank.

\subsection{Congruence and Quadratic Forms}

Recall from the preceding section that two real
symmetric matrices $A$ and $B$ are congruent if
\[
    B=S^TAS
\]
for some invertible matrix $S$.

Unlike similarity, congruence concerns the matrix
representation of a quadratic form rather than
that of a linear operator, and it preserves the rank of a matrix, though,
it does not, in general, preserve its eigenvalues.
The following example illustrates this distinction.

\begin{example}[Congruence Does Not Preserve Eigenvalues]

Consider the two real symmetric matrices
\[
    A=
    \begin{pmatrix}
        1&0\\
        0&-2
    \end{pmatrix},
    \qquad
    B=
    \begin{pmatrix}
        4&0\\
        0&-18
    \end{pmatrix}.
\]

We first observe that the eigenvalues of $A$ are
$1$ and $-2$, whereas those of $B$ are $4$ and
$-18$. Since similar matrices must have identical
eigenvalues, $A$ and $B$ cannot be similar.

However, are these matrices congruent?

To answer this question, recall that congruence
requires the existence of an invertible matrix
$S$ such that
\[
    B=S^TAS.
\]

Let us look for a diagonal matrix of the form
\[
    S=
    \begin{pmatrix}
        a&0\\
        0&b
    \end{pmatrix},
    \qquad ab\neq 0.
\]

Then
\[
\begin{aligned}
    S^TAS
    &=
    \begin{pmatrix}
        a&0\\
        0&b
    \end{pmatrix}
    \begin{pmatrix}
        1&0\\
        0&-2
    \end{pmatrix}
    \begin{pmatrix}
        a&0\\
        0&b
    \end{pmatrix}\\
    &=
    \begin{pmatrix}
        a^2&0\\
        0&-2b^2
    \end{pmatrix}.
\end{aligned}
\]

For this matrix to equal $B$, we require
\[
    a^2=4,
    \qquad
    -2b^2=-18.
\]

One possible choice is $a=2$ and $b=3$.
Therefore, taking
\[
    S=
    \begin{pmatrix}
        2&0\\
        0&3
    \end{pmatrix},
\]
we obtain
\[
    \boxed{S^TAS=B}.
\]

Furthermore,
\[
    \det(S)=6\neq 0,
\]
so $S$ is invertible. Hence, by the definition
of matrix congruence, $A$ and $B$ are congruent.

We have thus exhibited two matrices that are
congruent but not similar.

Interestingly, although their eigenvalues
differ, both matrices possess exactly one
positive and one negative eigenvalue.

This observation anticipates Sylvester's
Law of Inertia: congruence does not preserve
the individual eigenvalues, but it does
preserve the numbers of positive, negative,
and zero eigenvalues.
\end{example}



\begin{remark}
When $S$ is orthogonal, we have $S^{-1}=S^T$.
In this special case, congruence and similarity
take the same form:
\[
    S^TAS=S^{-1}AS.
\]
Consequently, orthogonal congruence preserves
the eigenvalues. For general invertible matrices,
however, this need not be true.
\end{remark}

\subsection{Why Rank Alone Is Not Enough}

Since multiplication by an invertible matrix
does not change rank, congruent matrices
necessarily have the same rank:
\[
\begin{aligned}
    \operatorname{rank}(S^TAS)
    &=\operatorname{rank}(AS)\\
    &=\operatorname{rank}(A).
\end{aligned}
\]

But is equality of ranks sufficient to
guarantee congruence?

Consider
\[
    A=
    \begin{pmatrix}
        1&0\\
        0&0
    \end{pmatrix},
    \qquad
    B=
    \begin{pmatrix}
        -1&0\\
        0&0
    \end{pmatrix}.
\]

Both matrices have rank one. Nevertheless,
they are not congruent.

To see why, observe that
\[
    x^TAx=x_1^2\geq 0
\]
for every $x\in\mathbb{R}^2$, whereas
\[
    x^TBx=-x_1^2\leq 0.
\]

Indeed, the quadratic form associated with $B$
takes strictly negative values, while the one
associated with $A$ never does.

If $B=S^TAS$ for some invertible $S$, then
\[
    x^TBx=(Sx)^TA(Sx)\geq 0
\]
for every $x$, contradicting the preceding
observation.

Therefore, equal rank is necessary but not
sufficient for congruence.

The obstruction in this example arises from
the signs of the quadratic forms, or equivalently,
from the signs of their eigenvalues.

This motivates the central question of our note:
\emph{What information about the eigenvalues
of a real symmetric matrix remains invariant
under congruence?}

The answer leads naturally to the concept of
\emph{inertia} and, ultimately, to Sylvester's
Law of Inertia.

\section{Inertia, Rank, Index, and Signature}

In the preceding section, we observed that congruent
matrices need not have identical eigenvalues. Nevertheless,
the signs of their eigenvalues appear to carry important
information about the underlying quadratic form.

This observation motivates the concept of \emph{inertia},
which plays a central role in the classification of real
symmetric matrices.

\subsection{Definition of Inertia}

Let $A\in\mathbb{R}^{n\times n}$ be a real symmetric
matrix. By the spectral theorem, all eigenvalues of $A$
are real, and there exists an orthogonal matrix $Q$
such that
\[
    Q^TAQ=
    \operatorname{diag}(\lambda_1,\ldots,\lambda_n).
\]

We may arrange these eigenvalues so that the positive
eigenvalues appear first, followed by the negative
eigenvalues, and finally the zero eigenvalues.

\begin{definition}[Inertia of a Symmetric Matrix]
Let $n_+$, $n_-$, and $n_0$ denote the numbers of
positive, negative, and zero eigenvalues of $A$,
respectively, counted with their multiplicities.

The \emph{inertia} of $A$ is the ordered triple
\[
    \boxed{\operatorname{In}(A)=(n_+,n_-,n_0)}.
\]
\end{definition}

Since every eigenvalue belongs to exactly one of
these three categories, we necessarily have
\[
    \boxed{n_++n_-+n_0=n}.
\]

The inertia therefore describes the distribution
of eigenvalues according to their signs, rather
than their actual numerical values.

\subsection{Rank, Index, and Signature}

Several important quantities can be expressed
directly in terms of inertia.

\begin{definition}[Rank, Index, and Signature]
For a real symmetric matrix $A$, we define:

\begin{enumerate}
    \item The \emph{rank} $r$ is the number of
    nonzero eigenvalues:
    \[
        r=n_++n_-.
    \]

    \item The \emph{index} $p$ is the number of
    positive eigenvalues:
    \[
        p=n_+.
    \]

    \item The \emph{signature} $s$ is the number
    of positive eigenvalues minus the number
    of negative eigenvalues:
    \[
        s=n_+-n_-.
    \]
\end{enumerate}
\end{definition}

Here, the term \emph{index} refers specifically
to the positive index of inertia. Some authors
use different terminology, but we shall follow
this convention throughout the note.

Since
\[
    r=n_++n_-
\]
and $p=n_+$, the number of negative eigenvalues
must be
\[
    n_-=r-p.
\]

Consequently, the signature becomes
\[
\begin{aligned}
    s
    &=n_+-n_-\\
    &=p-(r-p)\\
    &=2p-r.
\end{aligned}
\]

Thus,
\[
    \boxed{s=2p-r}.
\]

Equivalently, if the rank and signature are known,
we can recover the numbers of positive and negative
eigenvalues:
\[
    \boxed{
    n_+=\frac{r+s}{2},
    \qquad
    n_-=\frac{r-s}{2}
    }.
\]

Furthermore, by the rank-nullity theorem,
\[
    n_0=n-r.
\]

We therefore obtain
\[
    \boxed{
    \operatorname{In}(A)
    =(p,r-p,n-r)
    }.
\]

This shows that, for a symmetric matrix of fixed
size, its rank and index completely determine its
inertia. The signature is not an independent
quantity.

\subsection{An Illustrative Example}

Consider the real symmetric matrix
\[
    A=
    \begin{pmatrix}
        5&0&0&0\\
        0&2&0&0\\
        0&0&-3&0\\
        0&0&0&0
    \end{pmatrix}.
\]

Since $A$ is already diagonal, its eigenvalues
are immediately obtained:
\[
    \lambda_1=5,\qquad
    \lambda_2=2,\qquad
    \lambda_3=-3,\qquad
    \lambda_4=0.
\]

There are two positive eigenvalues, one negative
eigenvalue, and one zero eigenvalue. Hence,
\[
    \boxed{\operatorname{In}(A)=(2,1,1)}.
\]

The corresponding rank, index, and signature are
\[
\begin{aligned}
    r&=2+1=3,\\
    p&=2,\\
    s&=2-1=1.
\end{aligned}
\]

As expected, the identity $s=2p-r$ gives
\[
    s=2(2)-3=1.
\]

\subsection{Geometric Interpretation}

The inertia of a symmetric matrix also describes
important geometric properties of its associated
quadratic form.

To see this, suppose that
\[
    A=QDQ^T,
\]
where $Q$ is orthogonal and
\[
    D=\operatorname{diag}
    (\lambda_1,\ldots,\lambda_n).
\]

Introducing the orthogonal coordinates
$y=Q^Tx$, we obtain
\[
\begin{aligned}
    q(x)
    &=x^TAx\\
    &=x^TQDQ^Tx\\
    &=y^TDy\\
    &=\sum_{i=1}^{n}\lambda_i y_i^2.
\end{aligned}
\]

Thus, in an orthonormal eigenbasis, every
quadratic form becomes a sum of squared
coordinates weighted by its eigenvalues.

Positive eigenvalues contribute positively
to the quadratic form, negative eigenvalues
contribute negatively, and zero eigenvalues
make no contribution.

For instance, the matrix in our previous
example defines
\[
    q(x)=5x_1^2+2x_2^2-3x_3^2.
\]

The form is strictly positive for every
nonzero vector in the subspace
\[
    \operatorname{span}\{e_1,e_2\},
\]
and strictly negative for every nonzero
vector in
\[
    \operatorname{span}\{e_3\}.
\]

Moreover, the vector $e_4$ belongs to the
null space of $A$, so
\[
    Ae_4=0.
\]

These subspaces have dimensions two, one,
and one, respectively, precisely matching
the inertia $(2,1,1)$.

\begin{remark}
A vector satisfying $q(x)=0$ need not belong
to the null space of $A$. For example, if
$A=\operatorname{diag}(1,-1)$, then
$q(1,1)=0$, although $A(1,1)^T\neq 0$.
The zero component of inertia counts zero
eigenvalues, or equivalently, the dimension
of $\ker(A)$.
\end{remark}

We are now ready to formulate the central result of
this note: Sylvester's Law of Inertia. Recall that
the inertia of a real symmetric matrix is the ordered
triple $(n_+,n_-,n_0)$, recording the numbers of
positive, negative, and zero eigenvalues, respectively.
Sylvester's Law of Inertia asserts that these three
numbers remain unchanged under congruence, even
though the eigenvalues themselves may change.

More remarkably, the converse is also true: two
real symmetric matrices of the same size are
congruent if they have identical inertia.
Thus, inertia provides a complete classification
of real symmetric matrices under congruence.
We shall prove both directions of this remarkable
result in the next section.

% --------------------------------------------------
% Sylvester's Law of Inertia
% --------------------------------------------------

\section{Sylvester's Law of Inertia: Statement and Proof}

As discussed in the preceding
sections, congruence does not generally preserve
the individual eigenvalues of a symmetric matrix.
Nevertheless, it preserves the numbers of positive,
negative, and zero eigenvalues. Sylvester's Law of Inertia states that these numbers
completely characterize the congruence classes of
real symmetric matrices.

\subsection{Statement of the Theorem}

\begin{theorem}[Sylvester's Law of Inertia]
Let $B,C\in\mathbb{R}^{n\times n}$ be real symmetric
matrices. Then $B$ and $C$ are congruent if and
only if they have the same inertia.

Equivalently,
\[
    \boxed{
    B\text{ and }C\text{ are congruent}
    \iff
    \operatorname{In}(B)=\operatorname{In}(C).
    }
\]

That is, there exists an invertible matrix $S$
satisfying
\[
    C=S^TBS
\]
if and only if $B$ and $C$ have the same numbers
of positive, negative, and zero eigenvalues,
counting multiplicities.
\end{theorem}

We will prove the theorem in two directions.

First, assuming that $B$ and $C$ are congruent,
we will prove that their inertia must agree.
This is the more challenging direction, and
our proof will rely on the spectral theorem,
orthonormal eigenbases, and a carefully
constructed linear transformation.

We will then establish the converse by showing
that two symmetric matrices with the same
inertia can be transformed into the same
diagonal matrix through congruence.

% --------------------------------------------------
% First Direction
% --------------------------------------------------

\subsection{Congruence Implies Equal Inertia}

\begin{proof}[Proof of the Forward Implication]


\textbf{Step 1. Establish equality of ranks.}

Suppose that $B$ and $C$ are congruent.
By definition, there exists an invertible
matrix $S$ such that
\[
    \boxed{C=S^TBS}.
\]

Since $S$ is invertible, so is $S^T$.
Multiplication by an invertible matrix
does not change rank. Therefore,
\[
\begin{aligned}
    \operatorname{rank}(C)
    &=\operatorname{rank}(S^TBS)\\
    &=\operatorname{rank}(BS)\\
    &=\operatorname{rank}(B).
\end{aligned}
\]

Let their common rank be $r$. Thus,
\[
    \boxed{
    \operatorname{rank}(B)
    =\operatorname{rank}(C)=r.
    }
\]

Because both matrices are $n\times n$,
they also have the same nullity:
\[
    \dim\ker(B)=\dim\ker(C)=n-r.
\]

Consequently, $B$ and $C$ have the same
number of zero eigenvalues.

Our remaining task is to prove that they
have the same numbers of positive and
negative eigenvalues.

\medskip

\textbf{Step 2. Assume that the indices differ.}

Let $p$ denote the number of positive
eigenvalues of $B$, and let $q$ denote
the number of positive eigenvalues of $C$.

Suppose, for contradiction, that
\[
    p\neq q.
\]

Since the roles of $B$ and $C$ can be
interchanged, we may assume without
loss of generality that
\[
    \boxed{p<q}.
\]

Indeed, congruence is symmetric: if
$C=S^TBS$, then
\[
    B=S^{-T}CS^{-1}.
\]

Since both matrices have rank $r$,
the number of negative eigenvalues of
$B$ is $r-p$, whereas that of $C$
is $r-q$.

Our assumption $p<q$ therefore implies
\[
    r-p>r-q.
\]

Thus, under our assumption, $B$ has
fewer positive eigenvalues but more
negative eigenvalues than $C$.

We shall show that this is impossible
if the matrices are congruent.

\medskip

\textbf{Step 3. Choose orthonormal eigenbases.}

Since $B$ is real symmetric, the spectral
theorem guarantees the existence of an
orthonormal eigenbasis
\[
    \{v_1,v_2,\ldots,v_n\}
\]
such that
\[
    Bv_i=a_iv_i,
\]
where $a_i$ is the eigenvalue corresponding
to $v_i$.

We arrange the eigenvectors according to
the signs of their eigenvalues:
\[
\begin{cases}
    a_i>0, &1\leq i\leq p,\\
    a_i<0, &p+1\leq i\leq r,\\
    a_i=0, &r+1\leq i\leq n.
\end{cases}
\]

Likewise, since $C$ is real symmetric,
there exists an orthonormal eigenbasis
\[
    \{w_1,w_2,\ldots,w_n\}
\]
such that
\[
    Cw_i=e_iw_i,
\]
where
\[
\begin{cases}
    e_i>0, &1\leq i\leq q,\\
    e_i<0, &q+1\leq i\leq r,\\
    e_i=0, &r+1\leq i\leq n.
\end{cases}
\]

In particular,
\[
    \ker(B)
    =\operatorname{span}
    \{v_{r+1},\ldots,v_n\},
\]
and
\[
    \ker(C)
    =\operatorname{span}
    \{w_{r+1},\ldots,w_n\}.
\]

These relations follow because the
eigenvectors corresponding to zero
eigenvalues form an orthonormal basis
for the null space of a symmetric matrix.

Before proceeding, recall an important
consequence of congruence.

Since
\[
    C=S^TBS,
\]
for any $x\in\mathbb{R}^n$, define
\[
    y=S^{-1}x.
\]

Then $x=Sy$, and hence
\[
\begin{aligned}
    y^TCy
    &=y^TS^TBSy\\
    &=(Sy)^TB(Sy)\\
    &=x^TBx.
\end{aligned}
\]

Therefore,
\[
    \boxed{x^TBx=y^TCy,\qquad x=Sy.}
\]

This equality expresses the fact that
congruent matrices represent the same
quadratic form in different coordinates.

Notice that we must compare $x^TBx$
with $y^TCy$, where $x=Sy$, rather
than automatically using identical
coordinate vectors for both forms.

\medskip

\textbf{Step 4. Construct a linear transformation.}

We now introduce the central construction
of the proof.

Define the linear transformation
\[
    L:\mathbb{R}^n
    \longrightarrow
    \mathbb{R}^{p+(r-q)}
\]
by
\[
    \boxed{
    L(x)=
    \begin{pmatrix}
        v_1^TBx\\
        \vdots\\
        v_p^TBx\\
        w_{q+1}^TCS^{-1}x\\
        \vdots\\
        w_r^TCS^{-1}x
    \end{pmatrix}.
    }
\]


Let us examine this construction carefully.

The first $p$ components correspond to
the eigenvectors associated with the
positive eigenvalues of $B$.

The remaining $r-q$ components correspond
to the eigenvectors associated with the
negative eigenvalues of $C$.

Every component is linear in $x$.
Therefore, $L$ is a linear transformation.

Our objective is to find a vector
$x_0$ satisfying
\[
    L(x_0)=0,
\]
while ensuring that
\[
    Bx_0\neq 0.
\]

Such a vector will allow us to eliminate
the positive-eigenvalue contributions
of $B$ and the negative-eigenvalue
contributions of $C$.

The existence of this vector is guaranteed
by the rank-nullity theorem.

\medskip

\textbf{Step 5. Apply the rank-nullity theorem.}

Since
\[
    L:\mathbb{R}^n
    \longrightarrow
    \mathbb{R}^{p+(r-q)},
\]
we necessarily have
\[
    \operatorname{rank}(L)\leq p+(r-q).
\]

By the rank-nullity theorem,
\[
    \dim\ker(L)
    +\operatorname{rank}(L)=n.
\]

Consequently,
\[
\begin{aligned}
    \dim\ker(L)
    &=n-\operatorname{rank}(L)\\
    &\geq n-\bigl(p+r-q\bigr)\\
    &=n-r+(q-p).
\end{aligned}
\]

But our assumption was $p<q$.
Therefore,
\[
    q-p>0.
\]

Hence,
\[
    \boxed{\dim\ker(L)>n-r}.
\]

Recall, however, that
\[
    \dim\ker(B)=n-r.
\]

Thus,
\[
    \boxed{
    \dim\ker(L)>\dim\ker(B).
    }
\]

This inequality is the key observation.

A vector space of dimension greater
than $n-r$ cannot be contained entirely
in a vector space of dimension $n-r$.

Therefore, $\ker(L)$ cannot be contained
in $\ker(B)$.

Consequently, there exists a vector
$x_0\in\mathbb{R}^n$ such that
\[
    \boxed{
    x_0\in\ker(L),
    \qquad
    x_0\notin\ker(B).
    }
\]

Equivalently,
\[
    \boxed{
    L(x_0)=0,
    \qquad
    Bx_0\neq 0.
    }
\]

Now define
\[
    \boxed{y_0=S^{-1}x_0}.
\]

Since $C=S^TBS$, we obtain
\[
\begin{aligned}
    Cy_0
    &=CS^{-1}x_0\\
    &=S^TBS S^{-1}x_0\\
    &=S^TBx_0.
\end{aligned}
\]

Because $S^T$ is invertible and
$Bx_0\neq 0$, we conclude that
\[
    \boxed{Cy_0\neq 0}.
\]

We have therefore constructed two
corresponding coordinate vectors,
$x_0$ and $y_0=S^{-1}x_0$, satisfying
\[
    L(x_0)=0,
    \qquad
    Bx_0\neq 0,
    \qquad
    Cy_0\neq 0.
\]

We shall now study the quadratic form
in these two coordinate systems.

\medskip

\textbf{Step 6. Express $x_0$ in the eigenbasis of $B$.}

Since $\{v_1,\ldots,v_n\}$ is an
orthonormal basis for $\mathbb{R}^n$,
we may write
\[
    x_0=\sum_{i=1}^{n}c_iv_i.
\]

Recall that $L(x_0)=0$.

By the definition of $L$, its first
$p$ components must vanish:
\[
    v_k^TBx_0=0,
    \qquad 1\leq k\leq p.
\]

Substituting the expansion of $x_0$,
we obtain
\[
\begin{aligned}
    0
    &=v_k^TBx_0\\
    &=v_k^TB
      \left(\sum_{i=1}^{n}c_iv_i\right)\\
    &=\sum_{i=1}^{n}c_iv_k^TBv_i.
\end{aligned}
\]

Since $Bv_i=a_iv_i$,
\[
    0=\sum_{i=1}^{n}c_ia_iv_k^Tv_i.
\]

But the eigenvectors form an
orthonormal basis, so
\[
    v_k^Tv_i=\delta_{ki},
\]
where $\delta_{ki}$ denotes the
Kronecker delta.

Therefore, all terms vanish except
the term corresponding to $i=k$.
We obtain
\[
    \boxed{c_ka_k=0}.
\]

Since $a_k>0$ for $1\leq k\leq p$,
we must have
\[
    \boxed{c_k=0,\qquad 1\leq k\leq p}.
\]

Thus,
\[
    \boxed{
    x_0=\sum_{i=p+1}^{n}c_iv_i.
    }
\]

This is precisely what we intended:
the vector $x_0$ has no components
in the positive-eigenvalue directions
of $B$.

Furthermore, because $Bx_0\neq 0$,
at least one coefficient $c_i$ with
$p+1\leq i\leq r$ must be nonzero.

Indeed,
\[
    Bx_0=\sum_{i=p+1}^{r}c_ia_iv_i,
\]
and this expression cannot vanish
unless all coefficients in the indicated
range are zero.

\medskip

\textbf{Step 7. Show that $x_0^TBx_0<0$.}

Using the expansion obtained above,
we compute
\[
\begin{aligned}
    x_0^TBx_0
    &=
    \left(\sum_{i=p+1}^{n}c_iv_i\right)^T
    B\left(\sum_{j=p+1}^{n}c_jv_j\right)\\
    &=
    \sum_{i=p+1}^{n}
    \sum_{j=p+1}^{n}
    c_ic_jv_i^TBv_j.
\end{aligned}
\]

Since $Bv_j=a_jv_j$,
\[
    x_0^TBx_0
    =
    \sum_{i=p+1}^{n}
    \sum_{j=p+1}^{n}
    c_ic_ja_jv_i^Tv_j.
\]

By orthonormality, $v_i^Tv_j=\delta_{ij}$.
Therefore,
\[
    x_0^TBx_0
    =\sum_{i=p+1}^{n}a_ic_i^2.
\]

Since $a_i=0$ for $i>r$, this reduces to
\[
    \boxed{
    x_0^TBx_0
    =\sum_{i=p+1}^{r}a_ic_i^2.
    }
\]

Now observe that
\[
    a_i<0,
    \qquad p+1\leq i\leq r.
\]

Also, at least one coefficient
$c_i$ in this range is nonzero,
because $Bx_0\neq 0$.

Consequently, every summand is
nonpositive, and at least one
summand is strictly negative.

Hence,
\[
    \boxed{x_0^TBx_0<0}.
    \tag{1}
\]

Thus, the quadratic form represented
by $B$ takes a strictly negative
value at $x_0$.

\medskip

\textbf{Step 8. Express $y_0$ in the eigenbasis of $C$.}

Recall that
\[
    y_0=S^{-1}x_0.
\]

Since $\{w_1,\ldots,w_n\}$ is an
orthonormal eigenbasis for $C$,
we may write
\[
    y_0=\sum_{i=1}^{n}d_iw_i.
\]

Again, since $L(x_0)=0$, the last
$r-q$ components of $L(x_0)$ vanish:
\[
    w_k^TCS^{-1}x_0=0,
    \qquad q+1\leq k\leq r.
\]

Since $y_0=S^{-1}x_0$, this becomes
\[
    w_k^TCy_0=0.
\]

Substituting the expansion of $y_0$,
we obtain
\[
\begin{aligned}
    0
    &=w_k^TCy_0\\
    &=w_k^TC
      \left(\sum_{i=1}^{n}d_iw_i\right)\\
    &=\sum_{i=1}^{n}d_iw_k^TCw_i.
\end{aligned}
\]

Since $Cw_i=e_iw_i$,
\[
    0=\sum_{i=1}^{n}d_ie_iw_k^Tw_i.
\]

By orthonormality,
\[
    w_k^Tw_i=\delta_{ki}.
\]

Therefore,
\[
    \boxed{d_ke_k=0}.
\]

But
\[
    e_k<0,
    \qquad q+1\leq k\leq r.
\]

Hence,
\[
    \boxed{
    d_k=0,\qquad q+1\leq k\leq r.
    }
\]

Consequently,
\[
    \boxed{
    y_0=
    \sum_{i=1}^{q}d_iw_i
    +
    \sum_{i=r+1}^{n}d_iw_i.
    }
\]

Thus, $y_0$ has no components in
the negative-eigenvalue directions
of $C$.

Furthermore, since $Cy_0\neq 0$,
at least one coefficient $d_i$
with $1\leq i\leq q$ must be nonzero.

\medskip

\textbf{Step 9. Show that $y_0^TCy_0>0$.}

Using the orthonormal expansion,
we have
\[
\begin{aligned}
    y_0^TCy_0
    &=
    \left(\sum_{i=1}^{n}d_iw_i\right)^T
    C\left(\sum_{j=1}^{n}d_jw_j\right)\\
    &=
    \sum_{i=1}^{n}
    \sum_{j=1}^{n}
    d_id_jw_i^TCw_j\\
    &=
    \sum_{i=1}^{n}
    \sum_{j=1}^{n}
    d_id_je_jw_i^Tw_j\\
    &=
    \sum_{i=1}^{n}e_id_i^2.
\end{aligned}
\]

But we already established that
\[
    d_i=0,\qquad q+1\leq i\leq r,
\]
and
\[
    e_i=0,\qquad r+1\leq i\leq n.
\]

Hence,
\[
    \boxed{
    y_0^TCy_0
    =\sum_{i=1}^{q}e_id_i^2.
    }
\]

For every index in this sum,
\[
    e_i>0.
\]

Moreover, at least one coefficient
$d_i$ with $1\leq i\leq q$ is
nonzero, because $Cy_0\neq 0$.

Therefore, every summand is
nonnegative, and at least one
is strictly positive.

We conclude that
\[
    \boxed{y_0^TCy_0>0}.
    \tag{2}
\]

Thus, the quadratic form represented
by $C$ takes a strictly positive
value at $y_0$.

\medskip

\textbf{Step 10. Obtain the contradiction.}

Recall that
\[
    C=S^TBS
\]
and
\[
    x_0=Sy_0.
\]

Consequently,
\[
\begin{aligned}
    y_0^TCy_0
    &=y_0^TS^TBSy_0\\
    &=(Sy_0)^TB(Sy_0)\\
    &=x_0^TBx_0.
\end{aligned}
\]

Therefore,
\[
    \boxed{y_0^TCy_0=x_0^TBx_0}.
\]

However, equations (1) and (2)
established that
\[
    x_0^TBx_0<0
\]
while
\[
    y_0^TCy_0>0.
\]

Combining these results,
we obtain
\[
    0<y_0^TCy_0=x_0^TBx_0<0,
\]
which is impossible.

We have therefore reached a contradiction.

Hence, our assumption that the indices
of $B$ and $C$ differ must be false.
It follows that
\[
    \boxed{p=q}.
\]

We already established that the
matrices have the same rank $r$.
Consequently, they have the same
number of negative eigenvalues:
\[
    r-p=r-q.
\]

They also have the same number
of zero eigenvalues:
\[
    n-r.
\]

Thus, the numbers of positive,
negative, and zero eigenvalues
are identical for both matrices.

Therefore,
\[
    \boxed{
    \operatorname{In}(B)=\operatorname{In}(C).
    }
\]

This completes the proof that
congruence implies equality of inertia.

\end{proof}

% --------------------------------------------------
% Second Direction
% --------------------------------------------------

\subsection{Equal Inertia Implies Congruence}

We now establish the converse of the theorem.

Suppose that two real symmetric matrices
have the same inertia. We must show
that there exists an invertible matrix
relating them through congruence.

The central idea is that every real symmetric
matrix can be transformed, by congruence,
into a diagonal matrix whose entries
belong to $\{1,-1,0\}$.

\begin{proof}[Proof of the Converse Implication]

\textbf{Step 1. Assume equality of inertia.}

Let $B,C\in\mathbb{R}^{n\times n}$ be
real symmetric matrices satisfying
\[
    \operatorname{In}(B)
    =\operatorname{In}(C).
\]

Suppose their common inertia is
\[
    (p,r-p,n-r).
\]

Thus, each matrix has $p$ positive
eigenvalues, $r-p$ negative eigenvalues,
and $n-r$ zero eigenvalues.

Our objective is to construct an invertible
matrix $S$ such that
\[
    C=S^TBS.
\]

\medskip

\textbf{Step 2. Diagonalize the matrix $B$.}

By the spectral theorem, there exists
an orthogonal matrix $U$ such that
\[
    U^TBU=
    \operatorname{diag}
    (a_1,\ldots,a_n),
\]
where the eigenvalues are ordered
according to their signs:
\[
\begin{cases}
    a_i>0, &1\leq i\leq p,\\
    a_i<0, &p+1\leq i\leq r,\\
    a_i=0, &r+1\leq i\leq n.
\end{cases}
\]

We now introduce an invertible diagonal
matrix that will normalize the nonzero
eigenvalues of $B$.

Define
\[
    D_B=
    \operatorname{diag}
    \left(
    \frac{1}{\sqrt{|a_1|}},
    \ldots,
    \frac{1}{\sqrt{|a_r|}},
    1,\ldots,1
    \right).
\]

The entries corresponding to zero
eigenvalues are chosen to be $1$,
rather than attempting to divide
by zero.

Since every diagonal entry of $D_B$
is nonzero, $D_B$ is invertible.

Now observe that
\[
\begin{aligned}
    D_B^TU^TBUD_B
    &=
    D_B^T
    \operatorname{diag}(a_1,\ldots,a_n)
    D_B\\
    &=
    \operatorname{diag}
    \left(
    \frac{a_1}{|a_1|},
    \ldots,
    \frac{a_r}{|a_r|},
    0,\ldots,0
    \right).
\end{aligned}
\]

But
\[
    \frac{a_i}{|a_i|}
    =
    \begin{cases}
        1, &a_i>0,\\
        -1, &a_i<0.
    \end{cases}
\]

Therefore,
\[
    D_B^TU^TBUD_B
    =
    \begin{pmatrix}
        I_p&0&0\\
        0&-I_{r-p}&0\\
        0&0&0_{n-r}
    \end{pmatrix}.
\]

Let
\[
    P=UD_B.
\]

Since both $U$ and $D_B$ are
invertible, $P$ is invertible.

Thus,
\[
    \boxed{
    P^TBP=J_{r,p},
    }
\]
where
\[
    \boxed{
    J_{r,p}=
    \begin{pmatrix}
        I_p&0&0\\
        0&-I_{r-p}&0\\
        0&0&0_{n-r}
    \end{pmatrix}.
    }
\]

\medskip

\textbf{Step 3. Apply the same construction to $C$.}

Since $C$ is also real symmetric,
there exists an orthogonal matrix $V$
such that
\[
    V^TCV
    =\operatorname{diag}(e_1,\ldots,e_n).
\]

Because $B$ and $C$ have the same
inertia, we may order the eigenvalues
of $C$ so that
\[
\begin{cases}
    e_i>0, &1\leq i\leq p,\\
    e_i<0, &p+1\leq i\leq r,\\
    e_i=0, &r+1\leq i\leq n.
\end{cases}
\]

Define
\[
    D_C=
    \operatorname{diag}
    \left(
    \frac{1}{\sqrt{|e_1|}},
    \ldots,
    \frac{1}{\sqrt{|e_r|}},
    1,\ldots,1
    \right).
\]

As before, $D_C$ is invertible.

Setting
\[
    Q=VD_C,
\]
we obtain
\[
    \boxed{
    Q^TCQ=J_{r,p}.
    }
\]

Therefore, both matrices are congruent
to the same diagonal matrix:
\[
    \boxed{
    P^TBP=Q^TCQ=J_{r,p}.
    }
\]

\medskip

\textbf{Step 4. Construct the required congruence matrix.}

From the preceding equality,
\[
    P^TBP=Q^TCQ.
\]

Multiplying on the left by $Q^{-T}$
and on the right by $Q^{-1}$, we obtain
\[
    Q^{-T}P^TBPQ^{-1}=C.
\]

Rearranging the factors,
\[
    C=(PQ^{-1})^TB(PQ^{-1}).
\]

Define
\[
    \boxed{S=PQ^{-1}}.
\]

Since both $P$ and $Q$ are invertible,
the matrix $S$ is also invertible.

Consequently,
\[
    \boxed{C=S^TBS}.
\]

Thus, $B$ and $C$ are congruent.

This completes the proof that equality
of inertia implies congruence.

\end{proof}

% --------------------------------------------------
% Conclusion
% --------------------------------------------------

\subsection{Conclusion of the Proof}

We have now established both directions
of Sylvester's Law of Inertia.

First, we proved that congruent real
symmetric matrices necessarily have
the same inertia.

Conversely, we demonstrated that two
real symmetric matrices having the same
inertia can each be transformed into
the same canonical diagonal matrix,
and are therefore congruent.

Hence,
\[
    \boxed{
    B\text{ and }C\text{ are congruent}
    \iff
    \operatorname{In}(B)=\operatorname{In}(C).
    }
\]

% --------------------------------------------------
% Applications in Mathematical Statistics
% --------------------------------------------------

\section{Applications in Mathematical Statistics}

Sylvester's Law of Inertia has important applications
in mathematical statistics, where symmetric matrices
arise naturally in the study of covariance,
multivariate distributions, parameter estimation,
and optimization.

In this section, we illustrate several connections
between matrix congruence and statistical theory.

\subsection{Covariance Matrices and Linear Transformations}

Let $X$ be an $n$-dimensional random vector with
mean vector $\mu=\mathbb{E}[X]$ and covariance matrix
\[
    \Sigma_X
    =\mathbb{E}\left[
    (X-\mu)(X-\mu)^T
    \right],
\]
assuming the required second moments exist.

Recall that every covariance matrix is symmetric
and positive semidefinite. Indeed, for any vector
$v\in\mathbb{R}^n$,
\[
\begin{aligned}
    v^T\Sigma_Xv
    &=\mathbb{E}\left[
      v^T(X-\mu)(X-\mu)^Tv
      \right]\\
    &=\mathbb{E}\left[
      \bigl(v^T(X-\mu)\bigr)^2
      \right]\\
    &\geq 0.
\end{aligned}
\]

Consequently, all eigenvalues of $\Sigma_X$
are nonnegative.

Now suppose we introduce a new random vector
through an invertible linear transformation
\[
    Y=S^TX,
\]
where $S\in\mathbb{R}^{n\times n}$ is invertible.

The mean of $Y$ is
\[
    \mu_Y=S^T\mu.
\]

Therefore, its covariance matrix is
\[
\begin{aligned}
    \Sigma_Y
    &=\mathbb{E}\left[
      (Y-\mu_Y)(Y-\mu_Y)^T
      \right]\\
    &=\mathbb{E}\left[
      S^T(X-\mu)(X-\mu)^TS
      \right]\\
    &=S^T\mathbb{E}\left[
      (X-\mu)(X-\mu)^T
      \right]S.
\end{aligned}
\]

Hence,
\[
    \boxed{\Sigma_Y=S^T\Sigma_XS}.
\]

This is precisely a congruence transformation.

By Sylvester's Law of Inertia,
\[
    \boxed{
    \operatorname{In}(\Sigma_Y)
    =\operatorname{In}(\Sigma_X).
    }
\]

Since covariance matrices are positive
semidefinite, they have no negative eigenvalues.
Thus, if $\operatorname{rank}(\Sigma_X)=r$,
then
\[
    \operatorname{In}(\Sigma_X)=(r,0,n-r).
\]

The transformed covariance matrix has
the same inertia:
\[
    \operatorname{In}(\Sigma_Y)=(r,0,n-r).
\]

Therefore, an invertible linear transformation
may change the individual variances,
covariances, and eigenvalues, but it cannot
change the rank of the covariance matrix
or the number of its zero eigenvalues.

\begin{example}[A Singular Covariance Matrix]

Consider the covariance matrix
\[
    \Sigma_X=
    \begin{pmatrix}
        1&1\\
        1&1
    \end{pmatrix}.
\]

Its eigenvalues are $2$ and $0$, so
\[
    \operatorname{In}(\Sigma_X)=(1,0,1).
\]

This covariance structure corresponds to
two random variables having identical
centered values almost surely.

Now choose the invertible matrix
\[
    S=
    \begin{pmatrix}
        1&1\\
        0&1
    \end{pmatrix}.
\]

Under the transformation $Y=S^TX$,
the covariance matrix becomes
\[
\begin{aligned}
    \Sigma_Y
    &=S^T\Sigma_XS\\
    &=
    \begin{pmatrix}
        1&0\\
        1&1
    \end{pmatrix}
    \begin{pmatrix}
        1&1\\
        1&1
    \end{pmatrix}
    \begin{pmatrix}
        1&1\\
        0&1
    \end{pmatrix}\\
    &=
    \begin{pmatrix}
        1&2\\
        2&4
    \end{pmatrix}.
\end{aligned}
\]

The eigenvalues of $\Sigma_Y$ are $5$ and $0$.

Although the nonzero eigenvalue has changed,
the inertia remains
\[
    \boxed{
    \operatorname{In}(\Sigma_Y)
    =\operatorname{In}(\Sigma_X)
    =(1,0,1).
    }
\]

Thus, the covariance matrices are congruent
but have different eigenvalues, exactly as
predicted by Sylvester's Law of Inertia.

\end{example}

\begin{remark}
In statistical applications, the rank of
a covariance matrix represents the number
of linearly independent directions having
nonzero variance.

An invertible change of coordinates cannot
create or eliminate such directions.
It can only change how they are represented.
\end{remark}


\subsection{Whitening and Mahalanobis Distance}

Another important application of congruence
appears in the whitening of multivariate data.

Suppose that a random vector $X$ has a
positive definite covariance matrix $\Sigma$.
By the spectral theorem,
\[
    \Sigma=Q\Lambda Q^T,
\]
where all diagonal entries of $\Lambda$
are strictly positive.

We may therefore define
\[
    \Sigma^{-1/2}
    =Q\Lambda^{-1/2}Q^T.
\]

Consider the transformed random vector
\[
    Y=\Sigma^{-1/2}(X-\mu).
\]

Its covariance matrix is
\[
\begin{aligned}
    \operatorname{Cov}(Y)
    &=\Sigma^{-1/2}\Sigma\Sigma^{-1/2}\\
    &=I.
\end{aligned}
\]

Hence,
\[
    \boxed{\operatorname{Cov}(Y)=I}.
\]

This transformation is called
\emph{whitening}. It transforms the covariance
matrix into the identity matrix.

For example, if
\[
    \Sigma=
    \begin{pmatrix}
        4&0\\
        0&9
    \end{pmatrix},
\]
then
\[
    S=
    \begin{pmatrix}
        \frac12&0\\
        0&\frac13
    \end{pmatrix}
\]
satisfies
\[
    \boxed{S^T\Sigma S=I}.
\]

Notice that the eigenvalues have changed
from $4$ and $9$ to $1$ and $1$.
Nevertheless, both matrices have inertia
$(2,0,0)$.

Sylvester's Law of Inertia explains why
this normalization is possible: every
real symmetric positive definite matrix
is congruent to the identity matrix.

An important consequence appears in the
Mahalanobis distance, defined for a
positive definite covariance matrix by
\[
    d_M^2(x,\mu)
    =(x-\mu)^T\Sigma^{-1}(x-\mu).
\]

Under an invertible change of coordinates,
\[
    y=S^T(x-\mu),
    \qquad
    \Sigma_Y=S^T\Sigma S,
\]
we obtain
\[
    \Sigma_Y^{-1}
    =S^{-1}\Sigma^{-1}S^{-T}.
\]

Consequently,
\[
\begin{aligned}
    y^T\Sigma_Y^{-1}y
    &=(x-\mu)^TS
      S^{-1}\Sigma^{-1}S^{-T}
      S^T(x-\mu)\\
    &=(x-\mu)^T\Sigma^{-1}(x-\mu).
\end{aligned}
\]

Thus, the Mahalanobis distance is invariant
under invertible linear changes of coordinates.

This illustrates a broader geometric principle:
although a covariance matrix changes with the
coordinate system, the statistical geometry
it describes can remain unchanged.


\subsection{Fisher Information and Reparameterization}

A further application arises in statistical
parameter estimation.

Let $\theta\in\mathbb{R}^n$ be a parameter
vector and let $\ell(\theta;Z)$ denote
the log-likelihood function associated
with an observation $Z$.

The score vector is defined by
\[
    g_\theta
    =\nabla_\theta\ell(\theta;Z).
\]

Under the usual regularity assumptions,
the Fisher information matrix can be written as
\[
    I_\theta
    =\mathbb{E}[g_\theta g_\theta^T].
\]

The Fisher information matrix is symmetric
and positive semidefinite, since
\[
    v^TI_\theta v
    =\mathbb{E}[(v^Tg_\theta)^2]\geq 0.
\]

Now suppose that we introduce a new parameter
vector $\phi$ through a smooth, invertible
reparameterization
\[
    \theta=h(\phi).
\]

Let
\[
    J=\frac{\partial\theta}{\partial\phi}
\]
denote its Jacobian matrix, assumed invertible
at the parameter value under consideration.

By the multivariable chain rule,
\[
    g_\phi=J^Tg_\theta.
\]

Consequently,
\[
\begin{aligned}
    I_\phi
    &=\mathbb{E}[g_\phi g_\phi^T]\\
    &=\mathbb{E}[J^Tg_\theta g_\theta^TJ]\\
    &=J^T\mathbb{E}[g_\theta g_\theta^T]J.
\end{aligned}
\]

Therefore,
\[
    \boxed{I_\phi=J^TI_\theta J}.
\]

Once again, we obtain a congruence
transformation.

By Sylvester's Law of Inertia,
\[
    \boxed{
    \operatorname{In}(I_\phi)
    =\operatorname{In}(I_\theta).
    }
\]

Thus, a smooth, locally invertible
reparameterization preserves the rank
and inertia of the Fisher information matrix.

In particular, if the Fisher information
matrix is positive definite, it remains
positive definite under such a
reparameterization.

Likewise, if the Fisher information matrix
is singular, an invertible reparameterization
cannot remove that singularity.

This is important in parameter estimation
because singular Fisher information can
indicate a lack of local identifiability:
certain parameter directions may not be
distinguishable from the available
statistical information.


\subsection{Likelihood Curvature and Critical Points}

The full significance of inertia becomes
especially apparent when studying the
curvature of statistical objective functions.

Suppose that $\ell(\theta)$ is a twice
continuously differentiable log-likelihood
function, and let
\[
    H_\theta
    =\nabla_\theta^2\ell(\theta)
\]
denote its Hessian matrix.

Since the Hessian is symmetric, its inertia
describes the numbers of positive, negative,
and zero curvature directions.

Now consider a smooth invertible
reparameterization
\[
    \theta=h(\phi),
\]
with Jacobian $J$.

At a stationary point, where
\[
    \nabla_\theta\ell(\theta)=0,
\]
the multivariable chain rule gives
\[
    \boxed{
    H_\phi=J^TH_\theta J.
    }
\]

Thus, the Hessians in the two parameterizations
are congruent at corresponding stationary points.

Sylvester's Law of Inertia therefore implies
\[
    \boxed{
    \operatorname{In}(H_\phi)
    =\operatorname{In}(H_\theta).
    }
\]

This has an important interpretation.

If the Hessian is negative definite, the
stationary point is a strict local maximum.
If it is positive definite, the stationary
point is a strict local minimum.
If it has both positive and negative
eigenvalues, the stationary point is
a saddle point.

Because inertia is preserved under
invertible reparameterization, these
nondegenerate curvature classifications
do not depend on the choice of parameter
coordinates.

\begin{remark}
For nonlinear reparameterizations, the Hessian
transformation generally contains additional
terms involving first derivatives of the
objective function. These terms vanish at
stationary points, which is why the
congruence relation above holds there.
\end{remark}

The preceding examples illustrate how
Sylvester's Law of Inertia provides a
unifying principle across several areas
of mathematical statistics.

Whether we are transforming covariance
matrices, whitening multivariate data,
reparameterizing statistical models, or
examining the curvature of likelihood
functions, congruence preserves fundamental
structural information that is independent
of the particular coordinate system.

% --------------------------------------------------
% Conclusion and Historical Note
% --------------------------------------------------

\section{Conclusion}

In this note, we have explored the relationship
between quadratic forms, matrix congruence, and
Sylvester's Law of Inertia. Beginning with the
distinction between similarity and congruence,
we introduced the concept of inertia and established
that the numbers of positive, negative, and zero
eigenvalues of a real symmetric matrix remain
invariant under congruence.

We then presented a detailed proof of Sylvester's
Law of Inertia in both directions, using the
spectral theorem, orthonormal eigenbases, and
the rank-nullity theorem. Finally, we explored
its applications in mathematical statistics,
including covariance transformations, whitening,
Fisher information, and likelihood optimization.
The central lesson is that although individual
eigenvalues may change under congruence, their
distribution into positive, negative, and zero
values remains unchanged.

It would be rather unfair to conclude this note
without mentioning the remarkable mathematician
whose name is associated with this theorem.
James Joseph Sylvester (1814--1897) was an English
mathematician who played a pioneering role in
the development of modern matrix theory and
helped establish matrix algebra as a mathematical
subject in its own right. He introduced the
term \emph{matrix} in 1850, and his mathematical
exchanges with Arthur Cayley contributed
significantly to the emergence of modern
linear algebra.

Interestingly, Sylvester's academic career
was almost as extraordinary as his mathematics.
Despite ranking second in the Cambridge
Mathematical Tripos of 1837, he was initially
denied a degree because his Jewish faith
prevented him from subscribing to the religious
requirements of the university.

His teaching career was hardly uneventful either.
In 1841, Sylvester accepted a professorship at
the University of Virginia, but his appointment
ended rather dramatically in 1842 following a
dispute involving student misconduct.

According to a famous anecdote, Sylvester once
became involved in a confrontation with a student
and struck him with a sword-cane. The student
reportedly collapsed, leaving Sylvester briefly
convinced that he had killed him!

More colorful versions of the story portray
Sylvester as an exceptionally strict teacher,
whose students were too frightened to ask
questions. Indeed, according to such tales,
the real danger was not asking a question,
but being asked one by Sylvester himself.
If the unfortunate student failed to answer,
the punishment was supposedly rather more
physical than a poor examination grade!

These stories even portray Sylvester as a
mathematician repeatedly losing teaching
positions because of his temperament and
having to seek fresh appointments elsewhere.
Although such embellishments are difficult
to substantiate historically, his turbulent
experience at Virginia and his subsequent
moves between England and America certainly
make for an unusually eventful academic career.

In 1876, Sylvester joined the newly established
Johns Hopkins University, where he helped
develop advanced mathematical education
and founded the American Journal of Mathematics.
Among the younger scholars he mentored
were Thomas Craig, Fabian Franklin, and
Christine Ladd, later known as Christine
Ladd-Franklin, who went on to distinguished
intellectual careers. Particularly remarkable
was his encouragement of Christine Ladd at
a time when women were largely excluded from
advanced mathematical education. His willingness
to support and mentor a female mathematician
was unusually progressive for his era and
reveals another, perhaps less eccentric,
side of Sylvester's personality.

One might therefore jokingly wonder whether
Sylvester's legendary teaching methods
actually worked. After all, several of his
students became accomplished mathematicians.
Perhaps a little fear of the professor was
the secret ingredient---who knows!

Of course, the more convincing explanation
is Sylvester's extraordinary mathematical
insight, intellectual enthusiasm, and
commitment to developing younger scholars.

Whatever the truth behind the more entertaining
stories, Sylvester's lasting contribution
lies in the mathematics he helped create.
His Law of Inertia remains a beautiful example
of how a seemingly simple observation about
the signs of eigenvalues can reveal a profound
and coordinate-independent mathematical
structure.

Perhaps that is the most fitting conclusion:
mathematicians, with all their eccentricities,
eventually pass into history, but their
mathematical ideas continue to educate
generations long after they are gone.

% --------------------------------------------------
% References
% --------------------------------------------------

\begin{thebibliography}{99}

\bibitem{schultz2010}
A. Schultz,
\emph{Congruence; Sylvester's Law of Inertia},
Math 416 Lecture Notes,
University of Illinois at Urbana-Champaign,
April 22, 2010.
\url{https://palmer.wellesley.edu/~aschultz/w10/math416/coursenotes_and_handouts/coursenotes_100422.pdf}

\bibitem{rao2000}
A. Ramachandra Rao and P. Bhimasankaram,
\emph{Linear Algebra},
2nd edition,
Texts and Readings in Mathematics, Vol. 19,
Hindustan Book Agency, 2000.

\bibitem{anderson2003}
T. W. Anderson,
\emph{An Introduction to Multivariate Statistical Analysis},
3rd edition,
Wiley-Interscience, 2003.

\bibitem{casella2002}
G. Casella and R. L. Berger,
\emph{Statistical Inference},
2nd edition,
Duxbury, 2002.

\bibitem{nocedal2006}
J. Nocedal and S. J. Wright,
\emph{Numerical Optimization},
2nd edition,
Springer, 2006.


\end{thebibliography}
\end{document}
